\section{The closed fine-cell cover of a primary region}
\label{sec:area_law_primary_fine_cell_cover}

The primary regions in \cite[Section~11]{OpenAI2026AreaLaw} have an exact
decomposition into whole closed non-belt fine cells. One primary identifier
is retained even when its region is disconnected.

% Source: 10-geometry.tex, prop:two-families, lines 212--218 and 299--323,
% revision adc7f1241b42e322a6451854ab7e4b4c146bf78a.
% This is a local decomposition, not the full two-family proposition.

\begin{theorem}[Primary region as a finite union of closed fine cells]
    \label{thm:area_law_primary_fine_cell_cover}
    \lean{TNLean.PEPS.AreaLaw.Geometry.primaryBirthRegion_eq_iUnion_nonbeltCell_closure}
    \leanok
    \uses{def:area_law_primary_regions,def:area_law_belt_cells,
        def:area_law_nonbelt_pitch_index}
    Fix an arbitrary common origin, finite endpoint set and nonnegative
    integer neighborhood width. Let $k,\ell,p$ be natural exponents with
    $\ell\le k$ and $\ell\le p$, and put $m=2^{p-\ell}$. Choose canonical
    residues $0\le a,b<m$. Denote the actual fine-layer index set by
    $F_{k,\ell}$ and its belt subset by $\mathcal B$.
    For each signed pitch index $J\in\mathbb Z^2$, the primary birth region
    satisfies
    \begin{align}
        \overline{D_k\cap U_J}
        =\bigcup_{\substack{z\in F_{k,\ell}\setminus\mathcal B \\j(z)=J}}
        \overline{Q_\ell(z)}.
    \end{align}
    The equality includes empty layers, non-retained pitch indices,
    disconnected regions and the case $p=\ell$.
\end{theorem}
\begin{proof}\leanok
    \uses{thm:area_law_nonbelt_primary_containment,thm:area_law_fine_layer_count}
    Each closed non-belt fine cell with shifted pitch index $J$ lies in
    the primary birth region, giving one inclusion. For the other,
    choose the actual half-open fine cell containing a point of
    $D_k\cap U_J$. In either coordinate, let $z$ be its integer index
    and $J$ the pitch index, and put $r=z-a-mJ$ for the corresponding
    residue shift $a$. The strict pitch inequalities and the half-open
    cell inequalities imply $1\le r<m$. Euclidean division therefore
    gives $\lfloor(z-a)/m\rfloor=J$, and the residue is different from the belt residue.
    Doing this in both coordinates places the chosen cell in the indicated
    finite union. That union is closed, so it also contains
    $\overline{D_k\cap U_J}$. If $p=\ell$, the inequalities $1\le r<m$
    have no solution and both sides are empty.
\end{proof}

The equality supplies actual fine-cell witnesses for later interface
arguments. A disconnected primary may have several isolated contacts with
another region; two arbitrary common points do not by themselves specify
a nondegenerate interface segment.
